% ABOUTME: Conclusion section of the IEEE report.
% ABOUTME: Summary of Phase 1 findings and the plan for sequence models in Phase 2.
\section{Conclusion and Phase 2 Plan}
\label{sec:conclusion}

We predicted the remaining useful life of FD001 engines with engineered
degradation features and three classical regressors, under a protocol that
splits by engine, refits feature statistics inside every fold and scores
the test set once. The selected XGBoost model reaches a test RMSE of 15.36
cycles overall and 7.03 cycles on engines within 50 cycles of failure.
Against labels clipped at 125, Random Forest and XGBoost reach 14.02 and
14.08, ahead of published Random Forest and gradient boosting results and
0.8 to 1.4 cycles behind the best published deep models. The engineered
features and the tree ensembles both improve on the linear models with raw
sensors, the two tree ensembles cannot be separated on the test set, and
both rely mainly on the health index and the cycle number.

Phase 2 replaces the fixed 20-cycle summaries with models that read the
sensor sequence directly: an LSTM, a temporal convolutional network or a
Transformer encoder on fixed-length windows of the raw sensors. It adds an
ablation study and a hybrid model that feeds the engineered features into
the sequence model, which tests whether the hand-made features still help
once the history is learned. The sequence models will be evaluated under
the protocol of this report: FD001, five-fold cross validation grouped by
engine, training RUL capped at 125, one prediction per test engine scored
once, and the same RMSE, score $S$ and early and late stage metrics, so
that their results are directly comparable with Table~\ref{tab:test-results}.
